\documentclass[11pt,a4paper]{article}
\usepackage{fontspec}
\setmainfont{Noto Serif}
\setsansfont{Noto Sans}
\usepackage{babel}
\babelprovide[main,import]{russian}
\usepackage{amsmath,amssymb,mathtools,booktabs,geometry,fancyhdr,microtype,xcolor,hyperref}
\geometry{left=21mm,right=21mm,top=21mm,bottom=21mm}
\definecolor{ink}{HTML}{173D57}
\hypersetup{colorlinks=true,urlcolor=ink,linkcolor=ink,pdftitle={BGS-F. Спектральный перенос через резольвенту}}
\pagestyle{fancy}\fancyhf{}\lhead{\small BGS-F}\rhead{\small Спектральный перенос через резольвенту}\cfoot{\thepage}
\setlength{\headheight}{15pt}\setlength{\parindent}{1em}\setlength{\parskip}{2.5pt}
\setlength{\emergencystretch}{3em}\tolerance=3000
\newcommand{\Tr}{\operatorname{Tr}}
\newcommand{\EE}{\mathbb E}
\newcommand{\R}{\mathbb R}
\newcommand{\Z}{\mathbb Z}
\newcommand{\C}{\mathbb C}
\newcommand{\si}{\mathrm{Sine}_2}
\newcommand{\eps}{\varepsilon}
\newcommand{\sinc}{\operatorname{sinc}}
\newcommand{\sect}[1]{\section*{\sffamily\color{ink}#1}}
\newcommand{\subt}[1]{\par\medskip\noindent\textbf{#1.}\quad}
\newcommand{\T}{\mathbb T}
\newcommand{\dist}{\operatorname{dist}}
\newcommand{\Ad}{\operatorname{Ad}}
\newcommand{\op}{\mathrm{op}}
\newcommand{\HS}{\mathrm{HS}}
\newcommand{\Ph}{\Phi}
\newcommand{\CUE}{\mathrm{CUE}}
\begin{document}
\noindent{\small\sffamily\color{ink} МАТЕМАТИЧЕСКАЯ ФИЗИКА\hfill 20 НОЯБРЯ 2026}
\begin{center}{\LARGE\sffamily Спектральный перенос через резольвенту}\\[4pt]
{\large Обычные следы и близкие пары на степенях семи}
\end{center}
Рассматриваются исходные \(A_N,M_N\), размерности \(N=7^k\) и \(U_N(t)=e^{itA_N}M_N\). Физической силе \(u_N\) соответствует \(t=Nu_N\). В исходной задаче \(u_N=10^{-80}\). Период линейного оператора с точностью до скалярной фазы равен \(m=3N/7\).

\subt{Фиксированная сила}
Для \(N=7^{60}\) и \(N=7^{61}\) при исходном \(t=N/10^{80}\) доказано
\[
 |\Tr U_N(t)^m|/N>0.9526,\quad
 |\Tr U_N(t)^{2m}|/N>0.8914,\quad
 \Delta_{2,N}(b_*)>0.054.
\]
Старый ограниченный локальный тест даёт разность ожиданий с \(\si\) больше \(0.75\) также в этих двух размерностях.

При \(N=7^{62}>2.48\cdot10^{52}\) получено отличие от \(\si\) для другого ограниченного непрерывного локального теста \(G\) на близкие пары:
\[
                \EE G(\Xi_N)-\EE G(\si)>0.019 .
\]
Для этого достаточно доказать, что у более \(10\%\) узлов чистой сетки находятся как минимум две фазы в микроскопическом интервале длины \(0.27\) в развёрнутых единицах. Тест и его точные константы приведены в разделе 8.

\subt{Предельные результаты при убывающей силе}
\begin{itemize}
 \item При \(u_N=o(N^{-3/2}/\log N)\) весь локальный процесс сходится к решётке шага \(7/3\) со случайным непрерывным сдвигом и кратностями \(2(1+\nu_7(B+j))\), где \(B\) равномерен в \(\mathbb Z_7\). На каждом окне времён \(|r|\le TN\) нормированные обычные следы возмущённого и линейного операторов различаются на \(o(1)\), равномерно по \(r\).
 \item При \(\limsup |u_N|N^{3/2}\log N\le6\) ограниченный тест на близкие пары исключает \(\si\). Это включает ненулевые постоянные на граничном логарифмическом масштабе.
 \item При \(u_N=o(N^{-5/3}/\log N)\) сохраняется в пределе и полный парный тест \(\Delta_{2,N}(b_*)\). Его предел лежит в \((0.08443,0.08656)\). Более общо, разность всех фиксированных парных тестов Шварца двух спектров стремится к нулю.
\end{itemize}
Оценка использует распределение вейлевых мод по разностям линейных фаз, две стороны резольвенты и сохранение числа корней внутри контура. Для \(7^{62}\) далёкие узлы дополнительно устраняются дополнением Шура. Конечные сертификаты относятся к неизменному \(10^{-80}\). Асимптотический предел при этой фиксированной силе остаётся открытым.
\newpage
\sect{1. Энергия между узлами линейного спектра}
Пусть \(\lambda_j=\lambda_0e^{2\pi i j/m}\), \(j\in\Z/m\Z\), --- узлы линейного спектра, \(E_j\) --- соответствующие проекторы. Их кратности равны
\begin{equation}
    d_j=2\sum_{a=0}^{k-1}\mathbf1_{7^a\mid j}+\mathbf1_{j=0}.
                                                            \tag{1}
\end{equation}
Это следует из точного сопряжения \(M_N\) с перестановкой \(x\mapsto\gamma x\) над \(\Z/7^k\Z\), с точностью до общей фазы. Здесь \(\gamma=3+2\rho\), \(\rho^2=2\), \(\rho\equiv3\bmod7\). На каждом слое аддитивного порядка \(7^a\) две орбиты длины \(3\cdot7^{a-1}\); нулевой элемент даёт дополнительную фазу в узле \(j=0\).

Рассмотрим энергию переходов на разности \(h\) в невозмущённом собственном базисе:
\[
                    W_h=\sum_{j\bmod m}\|E_j A_N E_{j+h}\|_{\HS}^2 .
\]
Норма Гильберта--Шмидта здесь всюду \emph{ненормированная}. По ортогональности секторов \(W_h\) есть квадрат нормы проекции \(A_N\) на собственное подпространство \(\Ad M_N\) с собственным значением \(e^{-2\pi i h/m}\). В ортонормированном после деления на \(\sqrt N\) базисе \(T_N(v)\), сопряжение действует перестановкой \(S\) на метках. На одной орбите \(\mathcal O\) длины \(L\), если данное собственное значение допускается, вклад в \(W_h\) равен \(N/L\), умноженному на квадрат модуля суммы начальных коэффициентов орбиты с фазами единичного модуля. Если собственное значение не допускается, вклад нулевой.

Для исходных 32 коэффициентов все метки примитивны по \(7\), кроме \(v=(0,\pm7)\), с коэффициентами \(9/20\). При \(k\ge3\) орбиты примитивных меток определяются целой нормой \(Q(a,b)=a^2-2b^2\) и квадратным классом \(a+3b\bmod7\). Действительно, обе расщеплённые координаты обратимы, первая движется умножением на \(\gamma\), а произведение координат задаётся \(Q\). Элементы \(\langle\gamma\rangle\) составляют все единицы с квадратным первым разрядом. Разброс целых норм носителя меньше \(343\). Длины орбит примитивных меток равны \(m\); две непримитивные метки находятся на различных орбитах длины \(m/7\). Точные суммы модулей по группам дают
\begin{equation}
   W_h\le \frac{N}{m}\left(c_0+c_1\mathbf1_{7\mid h}\right),
   \qquad c_0=\frac{1243}{160000},\qquad c_1=\frac{567}{200}.
                                                             \tag{2}
\end{equation}
Здесь \(c_0=\sum_{\mathcal O\text{ прим.}}(\sum_{v\in\mathcal O}|c_v|)^2\). Все коэффициенты одной группы имеют равный модуль, так что это точный рациональный подсчёт. Вторая константа \(c_1=7\cdot2(9/20)^2\). Тот же исходный ряд даёт \(\|A_N\|_{\op}\le a_0=593/500\).

Для проверки (1) можно также вывести \(\sum_j d_j^2=(8N-5)/3\) непосредственно из числа орбит \(S\) на всём вейлевом базисе. На слое аддитивного порядка \(7^a\) их число \((7^{2a}-7^{2a-2})/(3\cdot7^{a-1})=16\cdot7^{a-1}\), плюс неподвижная нулевая метка. Число орбит равно размерности коммутанта \(M_N\), то есть \(\sum_j d_j^2\).
\newpage
\sect{2. Подсчёт собственных значений у каждого узла}
Фиксируем \(0<R\le1\). Вокруг узла \(\lambda_j\) в комплексной плоскости строим окружность \(\Gamma_j\) радиуса \(R/m\). Они попарно не пересекаются. Для циклического смещения \(d\) с \(|d|\le(m-1)/2\) положим
\[
 w_0=R^{-1},\qquad
 w_d=\bigl(2m\sin(\pi |d|/m)-R\bigr)^{-1}\quad(d\ne0),\qquad
 Z_j=\sum_{i\bmod m}\sqrt{w_{i-j}}\,E_i.
\]
Для всякого \(z\in\Gamma_j\) верно \(|\lambda_i-z|^{-1}\le m w_{i-j}\). Обозначим
\(\Delta=U_N(t)-M_N=(e^{itA_N}-I)M_N\).
Вдоль пути \(M_N+s\Delta-zI\), \(0\le s\le1\), разложение через два множителя \((M_N-zI)^{-1/2}\) сводит обратимость к обратимости \(I+s B(z)\), где
\[
     B(z)=(M_N-zI)^{-1/2}\Delta(M_N-zI)^{-1/2}.
\]
Корни в диагональном базисе \(M_N\) можно выбирать произвольно в каждой точке. Имеем \(B(z)=D(z)(m Z_j\Delta Z_j)D(z)\) с диагональным сжатием \(D(z)\), которое коммутирует с проекторами. Если \(\|m Z_j\Delta Z_j\|_{\HS}<1\), путь обратим на всём контуре. Принцип аргумента в конечной размерности сохраняет число собственных значений внутри \(\Gamma_j\), считая алгебраические кратности; для \(U_N(t)\) это число точно \(d_j\).

\subt{Усреднение квадрата нормы}
Пусть \(u_d,v_d\) --- произвольные неотрицательные веса на \(\Z/m\Z\), а \(Z_j(u)=\sum_i\sqrt{u_{i-j}}E_i\). Из (2) получаем
\begin{equation}
\frac1m\sum_j \big\||t|m\,Z_j(u) A_N Z_j(v)\big\|_{\HS}^2
\le t^2 N\left[c_0\Big(\sum_d u_d\Big)\Big(\sum_d v_d\Big)
               +c_1\sum_{a=0}^6 u^{(a)}v^{(a)}\right].
                                                             \tag{3}
\end{equation}
Здесь \(u^{(a)}=\sum_{d\equiv a\bmod7}u_d\). Действительно, усреднение вносит \(\sum_d u_d v_{d+h}\) в коэффициент при \(W_h\) и оставляет общий множитель \(t^2 m\). Сумма по всем \(h\) даёт произведение полных сумм, а условие \(7\mid h\) даёт произведения внутри классов.

\subt{Нелинейный остаток}
Запишем \(\Delta=(itA_N+J)M_N\), где \(\|J\|_{\op}\le (t a_0)^2/2\) по интегральному остатку для унитарной экспоненты. Неравенство \(w_i w_l\le(w_i^2+w_l^2)/2\) и расход квадратов по каждой строке и столбцу дают
\begin{equation}
   \frac1m\sum_j\|m Z_j J Z_j\|_{\HS}^2
          \le m N\|J\|_{\op}^2\sum_d w_d^2.                     \tag{4}
\end{equation}
При этом \(M_N\) справа не меняет норм, поскольку коммутирует с \(Z_j\). Сложив квадратные корни (3),(4), получаем верхнюю границу квадратного корня средней нормы \(\|m Z_j\Delta Z_j\|_{\HS}\) в \(L^2(j)\). Квадрат этой границы ограничивает долю узлов, у которых \(\|m Z_j\Delta Z_j\|_{\HS}\ge1\).
\newpage
\sect{3. Из хороших контуров в спектральное сопоставление}
Назовём узел хорошим, если доказано, что \(U_N(t)\) имеет внутри \(\Gamma_j\) ровно \(d_j\) собственных значений. Пусть \(\beta\) --- верхняя граница доли остальных узлов. Сопоставим в каждой хорошей окружности эти \(d_j\) фаз с линейным узлом. Доля \(\epsilon\) несопоставленных уровней, одинаковая в двух спектрах, по (1) удовлетворяет
\begin{equation}
     \epsilon \le \frac67\sum_{a=0}^{k-1}\min(\beta,7^{-a})+\frac1N,
       \qquad \epsilon\le(2k+1)\frac{m}{N}\beta .               \tag{5}
\end{equation}
Первое неравенство подсчитывает кратные \(7^a\) среди плохих индексов; второе использует максимум кратности. Сумма в первом ограничена \(C\beta[1+\log(1/\beta)]\) при \(\beta>0\), а при \(\beta=0\) равна нулю.

Если \(\lambda\) лежит на единичной окружности внутри \(\Gamma_j\), то модуль её углового отклонения, в радианах, не превосходит
\[
                  2\arcsin(R/(2m))\le (101/100)R/m .
\]
Последнее следует из границы производной \(\arcsin\) на \([0,1/10]\), поскольку \(m\ge21\). Для \(r\) порядка \(N\) отсюда следует равномерная оценка
\[
        N^{-1}|\Tr U_N(t)^r-\Tr M_N^r|
                      \le (101/100)|r|R/m+2\epsilon .           \tag{6}
\]
На скалярных временах \(r=l m\) есть нижняя оценка в фазе линейного возврата:
\begin{equation}
      \Re\left(\overline{\lambda_0}^{\,l m}\Tr U_N(t)^{l m}/N\right)
                  \ge 1-\tfrac12((101/100)lR)^2-2\epsilon .     \tag{7}
\end{equation}
Используем \(\cos x\ge1-x^2/2\) для сопоставленных уровней и границу \(-1\) для остальных.

\subt{Оценки весов}
Для \(d\ge1\)
\begin{equation}
 0\le w_d-\frac1{2\pi d}
      \le \frac{3d}{5m^2}+\frac{R}{18d^2},\qquad
                      w_d\le (3d)^{-1}.                        \tag{8}
\end{equation}
Действительно, \(2m\sin(\pi d/m)\ge4d\), а разность от \(2\pi d\) не превосходит \(\pi^3d^3/(3m^2)\) по синусному ряду. В разности обратных величин знаменатель не меньше \(3d\cdot6d\). Поэтому \(\sum_d w_d=O(R^{-1}+\log N)\) и \(\sum_d w_d^2\le R^{-2}+4/9\), используя \(\sum_{d\ge1}d^{-2}<2\). Суммы внутри отдельных классов по \(7\) позволят улучшить конечные константы.
\newpage
\sect{4. Асимптотический предел при более сильном убывающем кике}
При фиксированном \(R>0\) из (3),(4),(8) получаем
\begin{equation}
             \beta\le C_R\left(t^2N(1+\log N)^2+t^4N^2\right).
                                                             \tag{9}
\end{equation}
Если \(t=o(N^{-1/2}/\log N)\), то \(\beta\to0\) и по (5) \(\epsilon\to0\). Оценка (6) даёт \(\limsup\) для \(|r|\le TN\) не больше \(C_T R\); затем \(R\downarrow0\). Это доказывает утверждение об обычных следах. В частности, на времени \(3N/7\) нормированный след стремится по модулю к \(1\), как у линейного оператора.

\subt{Локальный закон}
Сравниваем точечные меры с одним и тем же равномерным началом развёртки. Для \(f\in C_c^1(\R)\) сдвиг сопоставленного уровня составляет \(O(R)\) в развёрнутых единицах. Усреднённая по началу абсолютная разность его вклада не превосходит \(C R\|f'\|_1/N\). Это следует интегрированием периодизованной функции и неравенством \(L^1\) для её сдвига. Каждый несопоставленный уровень даёт не более \(\|f\|_1/N\) в ожидании. Итак,
\[
       \EE_\alpha|\Xi_N^{U}(f)-\Xi_N^{M}(f)|
                             \le C R\|f'\|_1+2\epsilon\|f\|_1.
\]
Сначала \(k\to\infty\), затем \(R\downarrow0\) доказывают сходимость разности к нулю для каждого теста. Общий первый момент числа точек на компактах ограничен их длиной, потому что начало развёртки равномерно. Счётное семейство таких гладких тестов и теснота дают перенос слабого предела всего процесса.

Из (1) линейный предел определяется следующим образом. Сетка имеет шаг \(N/m=7/3\). При равномерном начале её непрерывный сдвиг равномерен в интервале длины \(7/3\). Для каждого \(D\) индекс ближайшего узла независимо равномерен по \(7^D\) при \(k-1\ge D\). По мере \(D\to\infty\) получаем равномерный \(B\in\mathbb Z_7\). Кратности \(2(1+\nu_7(B+j))\) следуют из (1), так как в любом фиксированном окне вероятность увидеть кратность, зависящую от разрядов после \(D\), не больше числа узлов, умноженного на \(7^{-D}\). Вклад дополнительного одиночного уровня при \(j=0\) исчезает по вероятности \(O(1/N)\). Этот предел имеет кратные атомы и отличается от \(\si\).

Физическое условие \(t=Nu_N\) превращается в \(u_N=o(N^{-3/2}/\log N)\). Все оценки в этой части даны для исходного оператора кика \(A_N\). Убывает только его сила.
\newpage
\sect{5. Сходимость полного парного теста}
Для \(\psi\in\mathcal S(\R)\) обозначим \(\psi_N(x)=\sum_{l\in\Z}\psi(N(x+l))\) и
\[
 P_\psi(V)=N^{-1}\sum_{i\ne j}\psi_N(\theta_i-\theta_j)
\]
по собственным фазам \(V\), с учётом кратностей. Докажем оценку
\begin{equation}
   |P_\psi(U_N(t))-P_\psi(M_N)|
                      \le C_\psi\left(R+\epsilon(1+N|t|+k)\right)
                                                             \tag{10}
\end{equation}
для фиксированного \(R\le1\), \(|t|=o(1)\).

\subt{Сопоставленные пары}
Если оба уровня сопоставлены, то после развёртки разность их фаз меняется на \(O(R)\). По теореме о среднем вклад ограничен \(C R\), умноженным на периодизацию \(\sup_{|y|\le C}|\psi'(\cdot+y)|\) по чистой сетке. По (1) и \(\sum_jd_j^2=(8N-5)/3\), для любого остатка \(h\)
\[
       N^{-1}\sum_j d_j d_{j+h}\le N^{-1}\sum_j d_j^2<8/3 .
\]
Дискретная сумма мажоранты \(\psi'\) по сетке шага \(7/3\) конечна. Поэтому общий вклад разности по сопоставленным парам не больше \(C_\psi R\).

\subt{Пары с несопоставленным уровнем}
В любой короткой дуге непрерывный унитарный путь \(e^{isA_N}M_N\) переносит собственные фазы не дальше \(a_0|t|/(2\pi)\), с попарным сопоставлением. Это следует из локальной аналитической нумерации через эрмитов логарифм и формулы \(2\pi\,\dot\theta=\langle v,A_Nv\rangle\). В кластере берут собственные значения сжатия \(A_N\). Укладка из (1) в дуге из \(J\) соседних чистых узлов даёт \(\sum_{j\text{ в дуге}}d_j\le C(J+k)\): для каждого \(a\) там не более \(J7^{-a}+1\) кратных \(7^a\). Следовательно, любой единичный интервал после развёртки содержит не больше \(C(1+N|t|+k)\) уровней возмущённого спектра. Для чистой сетки действует тот же довод с \(t=0\). Быстрое убывание \(\psi\) даёт ту же границу с \(C_\psi\) для модуля суммы всех парных вкладов, закреплённых за одним уровнем. После умножения на долю несопоставленных уровней получаем (10).

Если \(t=o(N^{-2/3}/\log N)\), то по второму неравенству (5) и (9) величина \(\epsilon(1+N|t|+k)=o(1)\) при каждом фиксированном \(R\). Главный член оценки имеет порядок \(N^2|t|^3(\log N)^3\). После \(R\downarrow0\) разность в (10) стремится к нулю. При \(\widehat\psi=w=|\widehat b_*|^2\) формула Пуассона даёт \(P_\psi(V)=N^{-2}\sum_r w(r/N)(|\Tr V^r|^2-N)\). Поэтому разность исходных \(\Delta_{2,N}(b_*)\) двух спектров есть половина разности этих \(P_\psi\). Предел для линейных матриц вычислен в части об обычном следе и заключён в \((0.08443,0.08656)\). Это доказывает заявленную асимптотику парного теста при \(u_N=o(N^{-5/3}/\log N)\).
\newpage
\sect{6. Равномерные конечные границы весов}
Для \(D=(m-1)/2\) и \(7\le L\le D\) разобьём веса по \(|d|\le L\) и \(|d|>L\). Пусть \(S_a^{\rm n}\), \(S_a^{\rm f}\) --- их суммы при \(d\equiv a\bmod7\), с \(a\in\{0,\ldots,6\}\). Вес \(d=0\) входит в первую сумму с \(a=0\). Используем элементарное интегральное правило для гармонической прогрессии:
\[
\sum_{\substack{1\le d\le L\\d\equiv a\bmod7}}\frac1d
   \le \frac1a+\frac17\log(L/a)\quad(1\le a\le6,\ L\ge a).
\]
Для \(a=0\) имеем \(\frac17(1+\log(L/7))\). Для участка \(L<d\le D\) граница каждого класса равна \(1/L+\frac17\log(D/L)\), включая класс \(a=0\). При сложении по двум знакам для класса \(a\ne0\) выступает пара \(a,7-a\), где \(1/a+1/(7-a)\le7/6\) и \(a(7-a)\ge6\).

Для линейной поправки \(3d/(5m^2)\) из (8) сумма при \(|d|\le L\) не больше
\(C_{\rm n}=3L(L+1)/(5m^2)\). На всей сетке отдельно для любого класса по \(7\) она не больше \(C_{\rm f}=3/70+3/(5m)\), так как в каждом положительном прогрессивном классе при \(d\le D\) не более \(D/7+1\) членов. Квадратичная поправка \(R/(18d^2)\) даёт границу \(2R/9\) по всей сетке, \(2R/(9\cdot49)\) для класса \(a=0\), исключая ноль, и \(R/(9L)\) для дальнего участка. Следовательно,
\[
\begin{aligned}
 S_0^{\rm n}&\le R^{-1}+\frac{1+\log(L/7)}{7\pi}
                         +C_{\rm n}+\frac{2R}{9\cdot49},\\
 S_a^{\rm n}&\le \frac{\log L}{7\pi}+\frac{7/6}{2\pi}
                       -\frac{\log6}{14\pi}+C_{\rm n}+\frac{2R}{9}
                            &&(a\ne0),\\
 S_a^{\rm f}&\le \frac{\log(D/L)}{7\pi}+\frac1{\pi L}
                              +C_{\rm f}+\frac{R}{9L}
                            &&(0\le a\le6).
\end{aligned}                                                   \tag{11}
\]
Для полных сумм можно брать \(L=D\), подставляя \(C_{\rm f}\) вместо \(C_{\rm n}\). Все числа в дальнейших сертификатах проверяются рационально с помощью \(3.14<\pi<22/7\), \(1.9<\log7<1.946\) и \(\log6>1.7\).

\subt{Исходная сила при \(k=60,61\)}
Выбираем \(R=1/5\). Тогда \(\log D<120\), и полные суммы удовлетворяют \(S_0<10.47\), \(S_a<5.70\) при \(a\ne0\). Подстановка в (3) с \(u=v=w\) даёт квадратный корень левой части меньше \(0.063\) для обоих \(k\). Из (4) квадратный корень вклада остатка меньше \(1/50000\). Отсюда
\[
          \beta<0.0041,\qquad
          \epsilon\le (6/7)(3\cdot0.0041+1/294)+1/N.
\]
Формула (7) даёт нижние границы \(0.9526,0.8914\) для нормированных следов при \(m,2m\). Сохраняя эти два положительных слагаемых парного теста и вычитая весь CUE-вклад, ограниченный \(C_w+2w(0)/N\), получаем \(\Delta_{2,N}(b_*)>0.054\). Подставляются известные точные границы \(w(3/7)>0.06854,\ w(6/7)>0.01599,\ C_w<0.02067,\ w(0)<0.1141\). Для старого ограниченного теста из части об обычном следе имеем \(\EE G_{\rm old}(\Xi_N)\le1-|\Tr U_N^m|/N\), а \(\EE G_{\rm old}(\si)\ge1-2/\sqrt{105}\), откуда разность больше \(0.75\).
\newpage
\sect{7. Дополнение Шура при \(N=7^{62}\)}
Здесь \(t=N/10^{80}\), \(R=7/20\), \(L=\lceil 50 |t|m\rceil\). Точные целочисленные неравенства \(L^2\le N\), \(D^2\le L^2N\), \(L<D\), дают \(\log L,\log(D/L)<60.4\). Из (11) получаем
\[
     S^{\rm n}_0<5.57,\qquad S^{\rm n}_a<2.98\ (a\ne0),\qquad
     S^{\rm f}_a<2.80\ (0\le a\le6).
\]
Для каждого центра \(j\) разобьём пространство через проекторы на \(|i-j|\le L\) и на дополнение, используя циклический представитель. Обозначим \(Z^{\rm n}_j,Z^{\rm f}_j\) ограничения весов \(Z_j\) на соответствующие части и положим
\[
 x_j=|t|m\|Z^{\rm n}_j A_N Z^{\rm n}_j\|_{\HS},\qquad
 y_j=|t|m\|Z^{\rm n}_j A_N Z^{\rm f}_j\|_{\HS}.
\]
Мажоранта весов \(mZ_j\) действует независимо от \(z\in\Gamma_j\). По эрмитовости \(A_N\) обе внедиагональные нормы его первого порядка ограничены \(y_j\). Формула (3) даёт
\[
 \EE_j x_j^2\le t^2 N\left[c_0\left(\sum_a S_a^{\rm n}\right)^2
                                 +c_1\sum_a(S_a^{\rm n})^2\right]<0.375274,
\]
\[
 \EE_j y_j^2\le t^2 N\left[c_0\left(\sum_a S_a^{\rm n}\right)
              \left(\sum_b S_b^{\rm f}\right)+c_1\sum_a S_a^{\rm n}S_a^{\rm f}\right]
                          <0.292646 .
\]

Для полного блока \(B_{\rm ff}(z)\) используем операторное ограничение
\[
       \|B_{\rm ff}(z)\|_{\op}\le |t|a_0 m/(4L-R)<1/100.
\]
Это верно для каждого \(z\in\Gamma_j\) и следует из \(\|\Delta\|_{\op}\le |t|a_0\). Остаток второго порядка в любой части \(B(z)\) имеет операторную норму \(\le m(t a_0)^2/(2R)<1/500\). Поэтому после устранения дальнего блока из \(I+sB(z)\), \(0\le s\le1\), отклонение ближнего дополнения Шура от единицы не больше
\[
        x_j+1/500+\frac{(y_j+1/500)^2}{1-1/100}
                    \le x_j+(103/100)y_j^2+1/400 .
\]
Последнее следует из \((y+e)^2\le(1+1/100)y^2+(1+100)e^2\). Если правая часть меньше \(1\), контур сохраняет кратность. Используем \(x\le(5/8)x^2+2/5\) при \(x\ge0\). По Маркову доля центров, у которых это достаточное условие не выполняется, не превышает
\[
          \beta\le \frac{(5/8)\EE_j x_j^2+(103/100)\EE_j y_j^2}
                          {1-1/400-2/5}<0.8971.                \tag{12}
\]
Так сохраняется кратность более чем у \(10\%\) узлов исходной сетки. Все константы в (11)--(12) проверяются целочисленно и рационально в приложенном коде.
\newpage
\sect{8. Ограниченный локальный тест на близкие пары}
Пусть \(h\) --- непрерывная функция со значениями в \([0,1]\), равная \(1\) на \([0.1,1.9]\), нулевая вне \([0,2]\) и линейная на двух переходах. Пусть чётная непрерывная функция \(j_0\) равна \(1\) при \(|s|\le0.27\), нулевая при \(|s|\ge0.28\) и линейная между границами. Для локально конечной точечной меры с возможными кратностями определим
\[
 G(\xi)=\min\left(1,\ \frac12\left[
  \iint h(x)h(y)j_0(y-x)\,d\xi(x)d\xi(y)
       -\int h(x)^2 j_0(0)\,d\xi(x)\right]\right).
\]
В скобках подсчитываются упорядоченные пары разных индексов; при совпадающих положениях точек кратность учитывается. Функционал непрерывен для смутной топологии на точечных мерах: компактная непрерывная функция двух переменных приближается суммами произведений, а массы компактных областей при смутной сходимости равномерно ограничены подходящим непрерывным мажорантом. Он ограничен от нуля до единицы.

Хороший узел из (12) несёт не меньше двух фаз по (1). После развёртки они лежат на расстоянии от центра не больше
\((101/100) R\,N/(2\pi m)\). Значит, полный диаметр меньше \(0.27\). Если центр попадает в интервал \([0.1+0.27/2,\,1.9-0.27/2]\), обе фазы лежат на единичной площадке \(h\), а \(j_0\) даёт вес \(1\). Тогда \(G=1\). По равномерности сдвига мера этого события для каждого узла \((1.8-0.27)/N\). События по различным центрам не пересекаются, потому что их шаг \(7/3\) больше длины интервала. Итак,
\[
          \EE_\alpha G(\Xi_{7^{62}})\ge (1-\beta)\frac37(1.8-0.27)
                         >0.0675.
\]

У \(\si\) вторая факториальная плотность равна \(1-(\sin\pi s/(\pi s))^2\). Для \(0\le s\le0.28\) она не превосходит \(\pi^2s^2/3\), по \(\sin\pi s/(\pi s)\ge 1-(\pi s)^2/6>0\). Поэтому
\[
          \EE G(\si)\le 2\int_0^{0.28}\frac{\pi^2s^2}{3}\,ds
                         <0.0482 .
\]
В оценке использована опора \(h\) длины \(2\); интегрирование по положительным разностям уже учитывает деление на два. Полученный строгий разрыв превышает \(0.019\). Этот сертификат использует ширину \(0.28\) в единицах средней щели; он относится к фиксированному \(N=7^{62}\) и исходному \(10^{-80}\).

\subt{Проверяемость и предел применимости}
Код \texttt{python3 resolvent/check.py --verify} проверяет константы (2)--(12), исходную силу и оба вида конечных свидетелей. Строгие числа для веса \(w\) перепроверяются командой \texttt{python3 zero\_mode/weights.py --verify}. Равномерный контроль (9) использует малость \(t\sqrt N\log N\). При \(t=N/10^{80}\) эта величина неограниченно растёт, поэтому из этих оценок нельзя получить предельную статистику фиксированного нелинейного оператора.

\newpage
\sect{9. Ненулевые силы на граничном логарифмическом масштабе}
Пусть \(\limsup |t_N|\sqrt N\log N=c<\infty\). Для произвольного фиксированного \(R\in(0,1]\) в аргументе раздела 7 теперь берём \(L=\lfloor\sqrt N\rfloor\). Для каждого \(a\bmod7\) обе суммы весов имеют ведущий член
\[
             S_a^{\rm n}=\frac{\log N}{14\pi}+O_R(1),
       \qquad S_a^{\rm f}=\frac{\log N}{14\pi}+O_R(1).
\]
Действительно, поправки в (8) дают ограниченную сумму; гармонические прогрессии в (11) имеют тот же ведущий член и снизу по интегральному сравнению. Для ближней части \(\log L=\tfrac12\log N+O(1)\), для дальней \(\log(D/L)=\tfrac12\log N+O(1)\). Поэтому (3) даёт для двух средних одну и ту же предельную границу
\[
     \limsup\EE_j x_j^2,\ \limsup\EE_j y_j^2
                \le \mu=\frac{(c_0+c_1/7)c^2}{4\pi^2}.
\]
Точный дальний блок имеет норму \(O(1/\log N)\), а операторный остаток второго порядка \(O_R(1/(\log N)^2)\), равномерно на всех контурах. При любом \(\eta>0\), для больших \(N\) достаточно условия
\[
                 x_j+(1+\eta)y_j^2+\eta<1 .
\]
Это следует из дополнения Шура и квадратичного поглощения равномерного остатка в \(y_j\). Для \(a>1/4\) имеем \(x\le a x^2+1/(4a)\). Применяя Маркова, а затем \(\eta\downarrow0\), получаем
\[
     \limsup\beta\le
       \frac{a+1}{1-1/(4a)}\mu.
\]
При \(a=(1+\sqrt5)/4\) множитель равен \(\varphi^2\), где
\(\varphi=(1+\sqrt5)/2\). Поэтому для
\[
               c<c_*:=\frac{2\pi}{\varphi\sqrt{c_0+c_1/7}}
\]
положительная доля узлов сохраняет кратность при каждом фиксированном \(R\), с предельной нижней границей независимо от \(R\). Тест из раздела 8 с достаточно малой шириной близкой пары исключает \(\si\): для исходного спектра ожидание ограничено снизу положительным числом, а в \(\si\) оно убывает кубически с шириной.

В частности, \(c_*>6\). При \(c\le6\) выбираем \(R=1/10\), прежнюю \(h\) и заменяем две границы \(j_0\) на \(0.08,0.09\). Полный диаметр хорошего кластера меньше \(0.08\). С точными числами \(c_0+c_1/7=66043/160000\), \(\varphi^2<2.619\) получаем
\[
 \limsup\beta<0.987,\qquad
 \liminf_{k\to\infty}\bigl[\EE G(\Xi_{7^k})-\EE G(\si)\bigr]>0.008 .
\]
Условие \(c\le6\) эквивалентно \(\limsup |u_N|N^{3/2}\log N\le6\). Это предельное опровержение для указанной области убывающих сил.
\end{document}
